\documentclass[11pt,a4paper]{article}
\usepackage[T1]{fontenc}
\usepackage[utf8]{inputenc}
\usepackage{amsmath,amssymb,amsthm}
\usepackage[margin=1in]{geometry}
\usepackage[colorlinks=true,linkcolor=blue,urlcolor=blue]{hyperref}
\theoremstyle{plain}
\newtheorem{theorem}{Theorem}
\newtheorem{lemma}[theorem]{Lemma}
\newtheorem{proposition}[theorem]{Proposition}
\newtheorem{corollary}[theorem]{Corollary}
\theoremstyle{definition}
\newtheorem{remark}[theorem]{Remark}

\title{The empirical recurrence for OEIS A233100, proved by transfer matrix}
\author{Adrian Perez Fontelles\\ \small Independent researcher}
\date{1 September 2026}
\begin{document}
\maketitle

\begin{abstract}
OEIS A233100 counts arrays over $\{0,\dots,3\}$ under a neighbour condition together with
two symmetry-breaking clauses: the upper-left entry is $0$, and $1$ must appear before
$2$ in row major order. The last of these is not local --- it compares first occurrences
across the whole array --- but it becomes local once the state carries which of the two values
was seen first. With that, $a(n)$ is a walk count on $S=1728$ vertices and the empirical
recurrence of order $5$ contributed by R.~H.~Hardin follows from a finite exact
computation.
\end{abstract}

\noindent\small 2020 Mathematics Subject Classification. 05A15, 05B45, 15B36.\normalsize

\section{The sequence and the conjecture}

OEIS A233100 is ``Number of 3 X n 0..3 arrays with no element x(i,j) adjacent to value 3-x(i,j) horizontally, diagonally or antidiagonally, top left element zero, and 1 appearing before 2 in row major order.''. It has offset $1$ and begins
\[
10, 74, 941, 11486, 141566, 1742447, 21452183, 264098522,\ \dots
\]
The entry carries the following comment:
\begin{quote}\small
Empirical: a(n) = 13*a(n-1) - 5*a(n-2) - 47*a(n-3) + 52*a(n-4) - 12*a(n-5) for n>6.
\end{quote}
As of the ``Last modified'' line on the live entry (Oct 07 2018 14:35:01, revision 7) the statement is
still recorded as empirical, and nothing on the entry records it as settled either way.

Written out, the claim is
\begin{equation}\label{eq:conj}
a(n)\;=\;13\,a(n-1) - 5\,a(n-2) - 47\,a(n-3) + 52\,a(n-4) - 12\,a(n-5)
\end{equation}
for all $n>6$.

\section{The three clauses}

The arrays have $3$ rows and $L=n$ columns; write $x_{t,u}$ for the entry in
column $t$ and row $u$. The neighbour condition is
\begin{equation}\label{eq:loc}
x_{t+d_1,\,u+d_2}\ne 3-x_{t,u}\qquad\text{for every offset }(d_1,d_2)\in\mathcal N=\{(1,0),\ (1,1),\ (1,-1)\}
\end{equation}
inside the array. The relation ``$y=3-x$'' is symmetric, so each unordered adjacency need
only be tested once; the offsets above are the forward representatives, one per direction,
and every one of them has a nonnegative column shift. A single column of state is therefore
enough for \eqref{eq:loc}: the cells of a column are settled by that column and the next.

Because the walk runs along columns while row major order runs along
rows, the two orders do not agree, and a single flag would not know which of the two
first occurrences came first. The state therefore carries one flag for each of the $3$
rows: $0$ while neither $1$ nor $2$ has been seen in that row, $1$ if $1$
was seen there first, $2$ if $2$ was. An array is accepted when the first row whose
flag is nonzero has flag $1$ --- row major order compares rows first and, within a
row, the earlier column --- or when no flag is set at all, in which case nothing is out of
order.

\begin{lemma}\label{lem:walk}
Let $\mathcal V$ be the set of pairs (a column, a flag vector), $|\mathcal V|=S=1728$, let
$(r,f)\to(s,f')$ be an edge when \eqref{eq:loc} holds for the cells of $r$ with $s$ beneath
it and $f'$ is $f$ updated by $s$, let $\mathbf{v}$ mark the vertices whose column begins
with $0$ and whose flag is the one that column alone produces, and $\mathbf{u}$ those whose
column satisfies \eqref{eq:loc} with nothing beneath it and whose flag is accepted. Then
\[
a(n)\;=\;\,\mathbf{v}^{\!\top}N^{\,L-1}\mathbf{u},\qquad L=n.
\]
\end{lemma}

\begin{proof}
A walk of length $L-1$ is a sequence of $L$ columns, and its flag component is determined by
them, so walks correspond to arrays. Along the walk the cells of column $t$ are tested once,
at the step leaving it, with the correct neighbour column; the cells of the last column are
tested by $\mathbf{u}$ with no column beneath, which is how the array behaves at its edge.
The flag component is by construction the record of first occurrences among the columns laid
down so far, so $\mathbf{u}$ imposes exactly the ordering clause, and $\mathbf{v}$ imposes
the upper-left clause. The state carries a column, not a boundary, so $L$ columns make a walk
of length $L-1$.
\end{proof}

\section{The criterion, and the computation}

Let $q(t)=t^{5}-\sum_i c_i\,t^{5-i}$ be the characteristic polynomial of
\eqref{eq:conj} and $G=N$.

\begin{lemma}\label{lem:crit}
With $n_0$ the least index for which $L\ge1$, $o=1n_0+0-1$ and
$h_j=G^{j}N^{o}\mathbf{u}$, we have for $j\ge5$
\[
a(n_0+j)-\sum_i c_i\,a(n_0+j-i)=\,\mathbf{v}^{\!\top}G^{\,j-5}w,
\qquad w=h_{5}-\sum_i c_ih_{5-i} .
\]
Hence \eqref{eq:conj} holds from that point on if and only if
$\mathbf{v}^{\!\top}G^{m}w=0$ for all $m\ge0$, and $m<S$ suffices.
\end{lemma}

\begin{proof}
Substitute Lemma~\ref{lem:walk} and factor. By Cayley--Hamilton $G^{S}$ is an integer
combination of $I,G,\dots,G^{S-1}$, so every later value repeats that combination of earlier
ones.
\end{proof}

\begin{theorem}
The recurrence \eqref{eq:conj} holds for every $n>6$, which is the statement of
Section~1.
\end{theorem}

\begin{proof}
The digraph was built directly from \eqref{eq:loc} and the flag update. The vector $w$ was
formed by $5$ applications of $G$ in exact integer arithmetic and
$\mathbf{v}^{\!\top}G^{m}w$ evaluated for $m=0,1,\dots$, again exactly. Every one of
those integers vanishes from the index corresponding to $n=6+1$ onwards, and the run
continued until $S+1$ consecutive zeros had been seen.
\end{proof}

\section{Verification}

Three checks, each independent of the algebra above.

First, the walk counts of Lemma~\ref{lem:walk} were compared against the entry's DATA at
every one of the $17$ published indices, with no shift fitted: the number of columns is
$L=n$ from the entry's own name and the offset says which index carries the first
published term. The values agree exactly. The ordering clause is where a misreading would
show first, and it is the reason the flags are per row rather than a single value.

Second, the recurrence \eqref{eq:conj} was evaluated on those published terms in exact
integer arithmetic, with no matrices involved, and vanishes wherever it applies.

Third, the annihilation test of Lemma~\ref{lem:crit} was carried to the full
Cayley--Hamilton bound in exact integer arithmetic rather than to a sampled prefix.

\begin{thebibliography}{9}
\bibitem{oeis} The OEIS Foundation, \emph{The On-Line Encyclopedia of Integer
Sequences}, \url{https://oeis.org}, 2026. Sequence A233100.
\bibitem{stanley} R.~P.~Stanley, \emph{Enumerative Combinatorics}, Volume 1, 2nd ed.,
Cambridge University Press, 2011. (Transfer-matrix method, Section 4.7.)
\bibitem{fs} P.~Flajolet and R.~Sedgewick, \emph{Analytic Combinatorics}, Cambridge
University Press, 2009.
\bibitem{horn} R.~A.~Horn and C.~R.~Johnson, \emph{Matrix Analysis}, 2nd ed., Cambridge
University Press, 2013.
\end{thebibliography}

\end{document}
